\section{Buffered circuits, mixing, and rare pins}\label{sec:cuts}

We first obtain double-$B$ circuits from an occupied-site annulus
estimate.  These circuits give conditional comparisons in fixed geometric
buffers and loop-count estimates inside arbitrary simple contours.  We
then construct the plane law and prove multiplicative comparisons for
rare pins.  Each stopping cut leaves its strict interior unexamined.

\subsection{An annulus input and local circuits}

Let $\mathbb H=\Tri^*$ be the honeycomb lattice.  Its hexagonal faces are
identified with the height vertices of $\Tri$; face distance means graph
distance in $\Tri$.  In lattice units write
\[
 e_1=(1,0),\qquad e_2=(1/2,\sqrt3/2),\qquad
 |s e_1+t e_2|_g=\max\{|s|,|t|,|s+t|\}.
\]
On face centers this norm is triangular graph distance.  For a center
face $z$, let $\Gamma_n(z)$ be the outer honeycomb contour of the union
of closed faces at distance at most $n$ from $z$, and let $\Lambda_n(z)$
be the honeycomb domain strictly inside that contour, excluding its
vertices and edges.  We suppress $z$ when it is the origin, and round
nonintegral indices down.  The union of these closed faces, its contour,
and the real norm ball of radius $n$ differ in radius by a bounded
additive constant.

The \emph{graph inradius} about $z$ of a simple contour $L$ surrounding
$z$ is the largest integer $n$ for which $\Lambda_n(z)$ and $\Gamma_n(z)$
lie strictly inside $L$.  It differs by a universal additive constant
from the minimum graph distance from $z$ to a face incident to $L$.
Indeed the smaller filled face balls cannot cross $L$, while the next
incident face layer reaches it.  In particular an exterior face incident
to $L$ can be chosen with center $u$ satisfying $|u-z|_g=n+O(1)$.

Choose one bipartite class $Y$ of honeycomb vertices.  Its nearest-neighbor
graph is triangular: two $Y$ sites are adjacent when joined by two
honeycomb edges.  At edge weight one, a site is \emph{occupied plus} if
all three incident faces have $B=+1$.  A circuit of occupied sites is a
circuit in this triangular site graph.

We use the annulus circuit estimate of
\cite[Theorem~2.3, p.~9; proof in Appendix~B]{OpenAIWeighted2026},
specialized to edge weight one: in a regular face-ball domain of radius $2r$, with
$B=-1$ on its boundary face layer and $W$ free, an occupied-plus $Y$
circuit surrounding the radius-$r$ ball has probability at least a
constant $p>0$, uniformly for all sufficiently large $r$.  Boundary-cell
factors are excluded in that theorem.  At edge weight one their removal
leaves exactly the uniform coherent-pair law with the indicated
one-spin boundary condition.  Only this annulus theorem is used from
\cite{OpenAIWeighted2026}.

Fix $0<\eps<1/50$ for the local circuit constructions below.

\begin{lemma}[Local double circuits]\label{lem:local-circuit}
There are $R>3$, $p>0$, and $r_0<\infty$ with the following property.  In a constrained coherent-pair law, suppose
that the radius-$Rr$ face-ball patch about $z$, including the two layers
of its bounding contour, contains no prescribed spins or $E_0$ edges.
For $r\ge r_0$, conditional on arbitrary $B$ values outside that patch,
there is probability at least $p$ of a simple double-$B$ circuit of either
specified sign whose curve surrounds $z$ and lies at norm radii
\begin{equation}\label{eq:local-circuit-annulus}
                (1-\eps)r\le |\,\cdot-z|_g\le(2+\eps)r.
\end{equation}
The circuit encloses the real norm ball of radius $(1-\eps)r$.
For a sign-$t$ circuit, additional $B=t$ pins in the patch are allowed.
\end{lemma}
\begin{proof}
Consider sign plus first.  Fix $B=-1$ on both face layers of the
radius-$2r$ regular contour.  By \cref{lem:spin-fkg} this lowers every
increasing event.  The exact cut law identifies its interior with the
one-spin minus-boundary law in the cited annulus theorem.  Bounded
changes in the chosen contour and integer radius are absorbed by the
fixed slack in \eqref{eq:local-circuit-annulus}.

Each edge of the occupied $Y$ circuit lifts to two honeycomb edges.
Both are double-$B$ plus: the two occupied endpoints prescribe all
incident face spins.  The resulting closed walk stays a bounded distance
from the site circuit and has the same winding about the inner ball.
Split the walk at repeated vertices and discard zero-winding pieces;
one obtains a simple double circuit with that winding.  The $Y$ lattice
and the face-center lattice are translates with the same graph norm, so
all the displacement errors just used are bounded independently of $r$.
For sufficiently large $r$ the resulting circuit satisfies
\eqref{eq:local-circuit-annulus} and has the asserted enclosure.

Choose, for example, $R=4$, increasing the lower threshold on $r$ to
accommodate all face stars and contour layers.  The comparison takes
place inside the prescribed empty patch, so arbitrary exterior $B$
conditioning leaves it valid.  Global sign reversal proves the minus
case.  Adding pins of the requested sign can only increase the circuit
probability by \cref{lem:spin-fkg}.
\end{proof}

\begin{lemma}[A circuit under the opposite cut condition]\label{lem:worst-cut}
There are fixed $K<\infty$ and $c>0$ with the following property.
Inside a regular dual cut of radius $Kn$, fix $B=-t$ on the inner boundary
layer and leave $W$ free.  With probability at least $c$, there is a
double-$B$ circuit of sign $t$ between radii $n$ and $2n$.  The constants
are independent of $n$, for all sufficiently large $n$.
\end{lemma}
\begin{proof}
Take a polygonal curve at norm radius $3n/2$.  Choose $r$ to be a
sufficiently small fixed multiple of $n$, and cover this curve cyclically
by annuli as in \eqref{eq:local-circuit-annulus}.  Put their centers on
the curve, including its corners, with successive distances between
$1.4r$ and $1.6r$.  Integer rounding preserves these inequalities with
slack for large $n$.  Choose $r/n$ small enough that the annuli and their
radius-$Rr$ comparison patches lie strictly between radii $n$ and $2n$.

The circuits in neighboring annuli must intersect.  Their enclosed
inner balls overlap, and each inner ball protrudes outside the other's
outer ball, since $\eps<1/50$.  Disjoint Jordan curves have disjoint or
nested interiors, and these two geometric properties exclude both
possibilities.  At a common honeycomb vertex, two double paths of the
same sign join with all three incident face spins of that sign.
Choose intersection points for consecutive pairs and join them along
the intervening circuits.  The resulting closed walk stays $O(r)$ from
the center polygon and is homotopic to it in the complement of the
radius-$n$ ball.  Extracting a simple circuit with nonzero winding gives
the required circuit between radii $n$ and $2n$.

There are only a fixed number of small annuli.  Their events are all
increasing in the order toward $t$, so \cref{lem:spin-fkg,lem:local-circuit}
give a positive lower bound for their intersection.  A fixed sufficiently
large $K$ keeps every comparison patch away from the opposite boundary
condition.  The constants do not depend on $n$.
\end{proof}

For clarity, a \emph{buffer} below means an open neighborhood in which
the local annuli, together with their larger worst-condition cuts, fit
with positive clearance.  Thus an obstacle can lie in the hole of a
collar without lying in the buffers used to construct the collar circuit.
We distinguish these local buffers from the entire disk bounded by the
eventual circuit.

\begin{lemma}[Circuits in buffers]\label{lem:buffer-circuits}
Consider a fixed positive-width polygonal collar, or a fixed positive-width
tube following a polygonal arc.  Suppose it has a buffer disjoint from
$E_0$ and from all prescribed $B$ spins.  At sufficiently fine mesh,
conditional on arbitrary $B$ values away from the buffer, there is a
probability at least $c>0$ of a double-$B$ circuit around the collar's hole,
or a connected double-$B$ path following the tube, of either specified sign.
For a sign-$t$ event, arbitrary additional $B=t$ pins in the buffer are
allowed.  The constant depends on the fixed geometry, not on the mesh,
the number of pins elsewhere, or their values.

Finitely many separated such requirements have a uniformly positive joint
probability, even if different buffers request different signs.  For a
logarithmic stack of separated similar annular buffers between radii $r$
and $L$, the probability that no circuit of a prescribed sign occurs is
at most
\begin{equation}\label{eq:stack-failure}
 C(r/L)^\alpha
\end{equation}
for some $\alpha>0$.  In lattice units this holds uniformly for $L\gg r\geq1$,
after adjustment of $C$.
\end{lemma}

\begin{proof}
First consider a small regular annulus whose larger cut lies in the
buffer.  Delete any favorable pins in that cut.  This can only lower the
probability of the desired event in the order toward $t$.  Fix the
opposite sign on both sides of the larger deterministic cut.  FKG lowers
the desired probability again, and the exact cut law identifies the
remaining problem with \cref{lem:worst-cut}.  This gives a uniform local
lower bound under every allowed exterior conditioning.

We explain the geometric gluing, with all comparison patches inside the buffer.  Cover the center curve of the tube
or collar by finitely many small regular annuli.  Choose their common
inner radius much smaller than the tube width.  Consecutive centers may
be separated by, for example, $3/2$ times that radius in the graph-distance
norm.  Their inner balls intersect, and each inner ball protrudes beyond
the other's outer ball.  Two surrounding circuits in these annuli must
therefore intersect: disjoint Jordan curves would have disjoint interiors
or nested interiors, and the two geometric conditions exclude these
possibilities.  Their union is connected and follows the prescribed
tube.  Around a closed center curve, its outer boundary contains a simple
circuit enclosing the hole; the annuli are small enough that their union
does not fill that hole.  At a common dual vertex, double paths of the
same sign have all three adjacent spin sites of that sign, so the gluing
does not lose the double-path property.  Bounded lattice rounding is
absorbed by the positive geometric clearances.

All the small annulus events for one sign are increasing in the same
order.  FKG multiplies their local lower bounds, giving a positive
constant depending only on the finite cover.  For separated buffers,
condition successively on the configurations outside the next buffer;
the preceding lower bound applies under each such conditioning.  This
also permits different signs in different buffers.

Finally choose a fixed scale ratio large enough that the enlarged buffers
of successive annular slots are disjoint.  There are at least
$c_2\log(L/r)-C_2$ slots.  At each trial, conditional on the previously
examined slots, success has probability at least $c_3>0$.  Failure in all
slots has probability at most $(1-c_3)^{c_2\log(L/r)-C_2}$, which is
\eqref{eq:stack-failure}.  The finitely many bounded lattice scales can be
absorbed in its prefactor.
\end{proof}

\subsection{Loop estimates inside simple contours}

For a simple honeycomb contour $L$, let $\PP_L$ be the loop $O(2)$ law
strictly inside $L$, with edge weight one and empty exterior: loops use
neither edges nor vertices of $L$, and have weight $2^{\#\mathrm{loops}}$.
It is represented by coherent pairs on the interior faces with both
spins plus on the inner boundary layer.  Indeed no wall can meet a
contour vertex, since all interior faces incident to that vertex are
boundary faces.  Conversely a loop strictly inside $L$ encloses no
boundary face.  Given the loops, assigning independently to each loop
the color $B$ or $W$ that changes across it reconstructs a unique coherent
pair with these boundary signs.  Thus the colors are independent fair
choices conditional on the loop geometry.

The first-spin marginal under this flat condition still admits the
comparisons above.  Embed $L$ in a larger regular spin domain, fix $B=+$
on the exterior faces and the inner boundary layer, and join all those
sites by $W$-equality edges.  The exterior and boundary face chain are
connected, so one connected equality set suffices.  Requiring its single
$W$ label to be plus has probability exactly $1/2$ conditional on every
$B$, by \cref{lem:spin-fkg}; it does not change the $B$ marginal.  The
resulting inside pair law is precisely the flat law just described.

\begin{proposition}[Loop estimates]\label{prop:loop-estimates}
There are constants $c_0>0$ and $0<\rho<1$ such that, for all sufficiently
large integer $n$, the following statements hold.
\begin{enumerate}[label=(\roman*)]
\item If the interior of a simple contour $L$ contains $\Lambda_{2n}(z)$,
then $\PP_L$ gives probability at least $c_0$ to a loop contained in
$\Lambda_{2n}(z)$ and surrounding $\Lambda_n(z)$.
\item If $L$ has graph inradius $n$ about $z$, then
\begin{equation}\label{eq:upper-loop-estimate}
 \PP_L\bigl(\text{at least two loops surround }\Lambda_{\rho n}(z)\bigr)
 \le 1-c_0.
\end{equation}
\end{enumerate}
There is no aspect-ratio or boundary-regularity assumption on $L$.
\end{proposition}
\begin{proof}
For (i), use two separated polygonal collars at norm radii, for example,
$1.3n$ and $1.7n$.  Their comparison buffers fit strictly between radii
$n$ and $2n$ and avoid the fixed boundary layers and equality edges.
\Cref{lem:buffer-circuits} gives a uniformly positive probability of a
double-$B$ plus circuit in the inner collar and a double-$B$ minus circuit
in the outer collar.  Different signs cause no association issue: reveal
the outer collar and apply the local conditional bound in the disjoint
inner buffer.  A $B$ wall separates these opposite-sign circuits.  No
wall can cross either double circuit, since all three face spins at each
of its vertices agree.  At least one wall loop therefore surrounds the
inner circuit and lies inside the outer one, proving (i).

For (ii), put $z=0$.  If a honeycomb circuit surrounds $\Lambda_s$, its
inner incident faces contain a triangular face-center circuit surrounding
the real norm ball of radius $s-b$, for a universal $b$.  To see this,
follow the inner incident faces along the contour.  Successive centers
are equal or adjacent, and the resulting closed walk tracks the contour
within bounded distance.  Its edges lie on the inner side.  Splitting at
repeated vertices gives a simple circuit with the same nonzero winding
about the smaller ball.  The same observation applies to a double
circuit whose face spins have one fixed sign.

Let $E$ be the event in \eqref{eq:upper-loop-estimate}, and let $E'$ also
require that its outermost surrounding loop changes $W$ and its second
outermost changes $B$.  The independent coloring gives
\begin{equation}\label{eq:two-loop-coloring}
                         \PP_L(E')=\tfrac14\PP_L(E).
\end{equation}
On $E'$, the outer loop is double-$B$ plus: no surrounding $B$ wall lies
outside it.  The faces immediately inside the next loop have $B=-1$.
Consequently there is a simple minus face-center circuit enclosing the
real norm ball of radius $\rho n-b$.  Other loops cannot alter these
incident signs without intersecting one of the two selected loops.

We select an enclosing plus cut using only $B$.  Extend the flat spins
as plus outside $L$, and start with all exterior faces reached.  Traverse
face adjacency unless the crossed honeycomb edge, including its endpoints,
is strictly inside $L$ and is double-$B$ plus.  A deterministic queue
queries the $B$ values needed to decide each attempted traversal, and
retains the full query record.  Complete this reachable closure.  If the
center is unreached, let $U$ be its component in the ordinary face-adjacency
graph of all unreached faces.  Each face adjacent to $U$ but outside it
is reached, and each edge separating the two sets is an eligible
strictly interior double-plus edge.  The reached set is connected to the
exterior, so $U$ has no holes.  At a honeycomb vertex only three faces
meet; the boundary of $U$ therefore has degree two at every vertex and
is a single simple circuit $C$.

Every original boundary face is reached across an ineligible contour
edge.  Moreover, every queried face in $U$ was tested from a reached
neighbor across a blocked edge, and hence is a plus inner-layer face of
$C$.  There are no other queried sites or original boundary constraints
in $U$.  Each completion consistent with the record reproduces the same
search and the same cut.  Across $C$ both $B$ values are plus, so coherence
places no restriction on $W$ there.  Thus \cref{lem:double-cut} gives
exactly the plus-$B$, free-$W$ cut law in $U$, even after all unobserved
exterior variables are summed out.  On $E'$ the outer selected loop
blocks the search, and $C$ lies on or outside that loop and encloses the
minus face circuit.
The search does not condition on $E'$ or on any loop color.

For a fixed realized $C$, let $A_C$ be the event of a simple minus
face-center circuit inside $C$ enclosing the real norm ball of radius
$\rho n-b$.  This is a decreasing event.  In a sufficiently large
regular free spin domain, prescribing plus on both layers of this fixed
$C$ gives its cut law and decreases $A_C$, by
\cref{lem:spin-fkg,lem:double-cut}.  Its probability in the stopped
interior is therefore at most its probability in that free domain.
We now bound the latter probability uniformly in $C$.

Choose an exterior face incident to $L$ with center $u$ satisfying
$|u|_g=n+O(1)$.  Fix $M>10(R+2)$, with $R$ from
\cref{lem:local-circuit}, and choose an integer $m$ such that
$(1-p)^m\le1/8$.  About $u$ test the annuli of that lemma at scales
\[
       d_i=\left\lfloor\frac{n}{20M^i}\right\rfloor,
                  \qquad i=0,\ldots,m-1.
\]
Test from large to small, revealing the $B$ spins in each tested annulus
and its incident face stars.  The radius-$Rd_{i+1}$ comparison patch
lies strictly inside the hole of every previously tested annulus for
large $n$.  Hence each trial has conditional success probability at
least $p$, and the probability of no success is at most $1/8$.
Take the free comparison domain large enough to contain $L$ and all
these patches with their bounding layers.

On success, the incident-face construction gives a simple plus
face-center circuit around $u$, enclosing the real norm ball of radius
$(1-2\eps)d_i$ and contained in the ball of radius $(2+2\eps)d_i$.
Set
\[
                       \rho=1-\frac{1}{200M^{m-1}}.
\]
For all sufficiently large $n$, the enclosed ball about $u$ overlaps
the radius-$(\rho n-b)$ ball about zero with linear slack.  Indeed the
smallest tested radius is $n/(20M^{m-1})+O(1)$, whereas
$(1-\rho)n=n/(200M^{m-1})$.  The plus circuit does not surround zero,
since its containing ball has radius at most $(2+2\eps)n/20$ and is
centered at distance $n+O(1)$ from zero.  A minus circuit counted by $A_C$ lies inside $C$, and hence inside $L$;
it cannot surround the exterior face center $u$.
The two circuits consequently have overlapping interiors with neither
containing the other.  They must intersect.  In the planar triangular
face graph, such an intersection contains a common spin site, impossible
for a plus and a minus circuit.

It follows that $\PP_{\rm free}(A_C)\le1/8$ for every realized $C$.
On $E'$ the search succeeds and $A_C$ occurs.  Averaging over its full
records therefore gives $\PP_L(E')\le1/8$.  Equation
\eqref{eq:two-loop-coloring} yields $\PP_L(E)\le1/2$.
Taking $c_0$ to be the smaller of $1/2$ and the constant in (i) proves
the proposition.  All lower thresholds on $n$ are imposed after the
fixed constants $R,p,M,m,\rho$ have been chosen.
\end{proof}

In a general flat height domain or a bounded skeleton face we use
\cref{prop:loop-estimates}\textup{(ii)} only after exposing an enclosing
level loop.  Its strict interior has the empty-exterior law: another
loop cannot use its vertices, and the unexplored inside retains the
weight $2^{\#\mathrm{loops}}$.  The contraction $\rho<1$ above is in
graph radii, as required for the distance-shell argument.

\subsection{Plane comparison}

We record the coupling mechanism behind our use of these circuits.  Take
two pair laws whose domains contain a large ball, with no constraints
there.  On a slightly smaller deterministic dual cut, compare each $B$
marginal with the law having minus fixed on both boundary layers.
Exact cutting makes this lower interior law common to both comparisons.

Run the exterior exploration of \cref{lem:double-cut} in the lower law,
querying the same sites in the other two laws.  At every stage retain,
for each marginal, its conditional law given its own queried values.
If these values are ordered, conditional FKG orders the next site's plus
probabilities: fixing the outer minus layer and then lowering the queried
values can only decrease that probability.  A fresh uniform random
variable therefore samples all three next spins in their required order.
The next query depends only on the revealed record, so this recursive
conditional sampling preserves each marginal.

Stop when the lower configuration has a plus double circuit.  Its cut
sites are plus in all three configurations.  The exploration has queried
only their exterior and these cut sites, so their remaining conditional
interior laws coincide.  Fill that interior identically.  Sample $W$ only
after this $B$ exploration; its components are free across the cut and can
also be matched.  The construction allows arbitrary constraints outside
the large ball, including $E_0$ equalities there.

The same construction compares a law with one obtained by adding extreme
$B$ pins in a small region.  Explore outward from that region in the more
adverse configuration.  A double circuit of the added sign matches the
configurations on the far side.  This version will be used when filling
small gaps in a pin curve.

\begin{proposition}[Plane law and quantitative local comparison]
\label{prop:plane-mixing}
Free coherent pairs in expanding regular hexagons converge on cylinder
events to a law $\mu_\infty$.  This law is invariant under lattice
translations, lattice symmetries, global sign changes, and interchange
of the colors.  If the domain of a finite pair law $\mu^{E_0,\eta}$
contains a ball of radius $L$ and has no constraints in that ball, then
its restriction to a concentric
ball of radius $r$ differs from that of $\mu_\infty$ in total variation
by at most $C(r/L)^\alpha$, for $L\gg r\geq1$.

The associated plane height increments exist and have the same spatial
symmetries.  Every finite feasible prescription of $B$ pins together with
$W=+1$ on finitely many connected paths has positive plane probability.
\end{proposition}

\begin{proof}
Apply the preceding common-lower-law coupling and search for a plus circuit
in a stack between the inner ball and the deterministic outer cut.  Its
failure probability is bounded by \eqref{eq:stack-failure}.  On success
the full pair configurations agree in the inner ball.  This bounds the
total variation distance between any two sufficiently large free-volume
laws, and proves the Cauchy property on each cylinder.  It gives the same
bound with one constrained law and one expanding free volume; then let
the latter volume tend to the plane.  Constants do not depend on the
constraints outside the empty ball.

Translated or reflected exhausting volumes give the same limit by this
comparison, proving the spatial symmetries.  The finite laws have the
stated spin symmetries.  Formula \eqref{eq:height-increment} determines
increments on every fixed path, and its vanishing curl persists in the
limit.  It therefore defines the plane height modulo one additive
integer.

To check positivity of a finite pin event, enclose all its sites and paths
by a double-$B$ cut.  Such a cut has positive probability by
\cref{lem:buffer-circuits}.  In its finite interior, take $W$ everywhere
plus and choose $B$ with the prescribed values and the cut value.  This
is a coherent extension and has positive conditional counting probability.
The exact cut law proves the assertion.  Uniform finite-volume estimates
pass to this positive-probability plane conditioning.
\end{proof}

When height increments rather than the two spin labels are considered,
one may equivalently anchor a free height at zero.  The global Gray
rotation $(B,W)\mapsto(W,-B)$ preserves \eqref{eq:height-increment} and
cycles the four residues at the anchor.  Each anchored height consequently
has the same four representatives in the free pair law.  This observation
will let us apply counting inequalities directly to anchored height arrays.

\subsection{A comparison for separated rare pin events}

A family of \emph{separated obstacles} consists of finitely many supports,
each contained in a bounded simply connected polygonal neighborhood $U_i$,
with the closures $\overline U_i$ pairwise disjoint.  Each support has a
surrounding polygonal collar inside $U_i$, with positive clearance for all
its local buffers and larger comparison cuts.  Their filled interiors
also lie in $U_i$.  Thus an isolating cut encloses only its own obstacle;
disjoint annular collars alone would not impose this condition.
After rescaling, these neighborhoods, collars, and clearances are fixed.
For example, finitely many lattice paths converging uniformly to disjoint
compact simple arcs admit such neighborhoods: choose sufficiently thin
regular neighborhoods of the arcs, with endcaps, and approximate their
boundaries polygonally.  Their closures remain disjoint at fine mesh.  On obstacle $i$, an event $P_i$ prescribes arbitrary
$B$ values and requires $W=+1$ on at most $m_i$ connected subpaths.  There
is a fixed bound on $\sum_i m_i$, but no bound on the number of pinned
sites along a subpath.  We use the same notation $P_i$ for the event and
write $p_i=\mu_\infty(P_i)$ for its probability.

The following comparison treats one obstacle while retaining any pins
already imposed on it.  This retention allows its plus and minus pins to
be prescribed in successive batches.

\begin{lemma}[Relative comparison for one obstacle]
\label{lem:relative-pin-comparison}
Let $\mu$ and $\nu$ be finite pair laws with second-color equality
constraints and first-color pins.  Their domains contain the same filled
isolating neighborhood $U$, and their constraints inside $U$ agree:
these consist of retained $B$ pins $R$ and equality edges $E_0$ on one
obstacle.  Its surrounding collar and buffers in $U$ contain no
constraints.  Outside $U$ the two laws may have different constraints.
If $A$ prescribes one sign $t$ on further sites of that obstacle,
consistently with $R$, then
\[
 c\nu(A)\leq\mu(A)\leq c^{-1}\nu(A).
\]
Here $c>0$ depends only on the collar and its buffers, not on the
retained pins, equality edges, remote constraints, or the probability
of $A$.  The same assertion holds with the colors interchanged.
\end{lemma}

\begin{proof}
Fix sign $t$ on a deterministic double layer $\Gamma$ in the collar.
Let $q$ be the probability of $A$ in its interior law, retaining $R$
and the internal $E_0$ constraints.  Exact cutting makes $q$ the same
for $\mu$ and $\nu$.  FKG in the order toward $t$ gives, for either
law $\lambda\in\{\mu,\nu\}$,
\begin{equation}\label{eq:rare-pin-upper}
 \lambda(A)\leq\lambda(A\mid B=t\text{ on }\Gamma)=q.
\end{equation}
In a thinner collar strictly inside $\Gamma$, an exterior stopping search finds
a double circuit of sign $t$ with probability at least $c_1>0$, by
\cref{lem:buffer-circuits}.  Its constant is uniform over $R$ and all
remote constraints: the search buffers contain none of these pins or
$E_0$ edges.  Conditional on the search, the inner law is the corresponding
closer-cut law.  In the deterministic-$\Gamma$ law, conditioning on this closer
sign-$t$ cut can only increase $A$, while preserving $R$.  Its probability
of $A$ is therefore at least $q$.  Consequently
\begin{equation}\label{eq:rare-pin-sandwich}
 c_1q\leq\lambda(A)\leq q,
 \qquad \lambda\in\{\mu,\nu\}.
\end{equation}
Comparison through this common $q$ proves the lemma.  The retained pins
may be extremely rare; their probability never enters the inequalities.
The stopping exploration has not inspected the strict interior, which
is essential for the lower bound.
\end{proof}

\begin{proposition}[Products of separated pin probabilities]
\label{prop:pin-product}
For a separated family as above, there are constants $0<c\leq C<\infty$
such that
\begin{equation}\label{eq:pin-product}
 c\prod_i p_i\ \leq\ \mu_\infty\Bigl(\bigcap_iP_i\Bigr)
 \ \leq\ C\prod_i p_i.
\end{equation}
The constants depend on the collars, their buffers, the number of
obstacles, and $\sum_i m_i$.  They do not depend on the mesh, the number
or the values of the microscopic $B$ pins, or the probabilities $p_i$.
The same estimates hold in all sufficiently large finite free volumes.
\end{proposition}

\begin{proof}
First handle the equality requirements for $W$.  Temporarily
ignore all $B$ pins.  For a collection of $m$ connected paths, let $E$
denote constancy of $W$ along each path and let $L$ require all labels to
be plus.  Conditional on $B$, a consistent pattern of labels has the same
probability as the all-plus pattern, and an inconsistent pattern has
probability zero.  Summing over the at most $2^m$ patterns gives
\begin{equation}\label{eq:constancy-label-comparison}
 \PP(L)\leq\PP(E)\leq 2^m\PP(L).
\end{equation}
Apply \cref{lem:relative-pin-comparison} with the colors interchanged to the
successive all-plus requirements on the obstacles.  Each conditional
factor is comparable to the corresponding one-obstacle factor: the
conditioning concerns only other, separated obstacles.  The all-plus probability thus factors up to
constants.  Equation \eqref{eq:constancy-label-comparison}, for the joint
family and for its individual members, gives the same conclusion for
their constancy events.  Denote the resulting equality edges by $E_0$.

Now condition on all these equalities and impose the $B$ requirements
obstacle by obstacle.  On each obstacle first impose its plus pins, then
its minus pins.  At either stage compare its conditional factor with the
same stage in the one-obstacle equality problem using
\cref{lem:relative-pin-comparison}.  At the minus stage the plus pins already
imposed on that obstacle belong to $R$ and are retained in both cut laws.
All previous pins on other obstacles and all other equality paths lie
outside $U_i$, so the two laws have exactly the same retained constraints
inside $U_i$.  Multiplication compares the joint conditional $B$-pin
probability with the product of the one-obstacle conditional probabilities.

Finally restore all the required $W$-plus labels.  Conditional on $B$ and
$E_0$, their probability is between $2^{-\sum_i m_i}$ and one; the same
statement applies in each one-obstacle problem.  These bounded factors,
together with the preceding two comparisons, prove
\eqref{eq:pin-product}.  The entire proof takes place in finite volumes
with uniform constants.  The plane assertion follows from
\cref{prop:plane-mixing} and positivity of each finite pin event.
\end{proof}
